\section{Related work}\label{sec:related}

\paragraph{The Six Birds series.} \FoundI{} defines finite theory packages and the six structural
roles P1--P6 \citep{Tsiokos2026FoundI}, and \FoundII{} supplies admissibility and quotient discipline
\citep{Tsiokos2026FoundII}. \FoundIII{} gives a finite audited interaction calculus: source-of-truth
tags, promotion gates, top-down channels, and the theorem that an instrument cannot certify its own
soundness at its own level \citep{Tsiokos2026FoundIII}. \FoundIV{} is the static catalog of
fifty-two laws whose split-pair descent criterion and source-refinement repair are the starting
point here \citep{Tsiokos2026FoundIV}. \FoundVI{} returns to an external analyst and studies laws
about whole runs \citep{Tsiokos2026FoundVI}; nothing in this paper depends on it. The adequacy
residual and its chain rule \citep{Tsiokos2026Adequacy}, currencies as constraints and shadow
prices \citep{Tsiokos2026Currency}, hiddenness and exposure \citep{Tsiokos2026Hiddenness}, the
probe-enlargement problem \citep{Tsiokos2026NeedleKiller}, predictive quotients and route
transport \citep{Tsiokos2026Holonomy}, agents and viability kernels \citep{Tsiokos2026Agents},
institutions as strict extensions \citep{Tsiokos2026Institutions} and the reflexive audit loop
\citep{Tsiokos2026Reflexive} are used where the laws need them.

Five questions left open in those papers motivate the catalog. \FoundII{} deferred the case of an
instrument whose soundness and completeness are at issue and which is itself part of the system.
The reflexive development kept its audit loop on the analyst's side. The probe-enlargement problem
asked which probe to add under which budget. The currency development named attention as a price
and the self-regulating accounting cell but did not develop them. The institutions development
asked for a carrier that acquires, compiles and reactivates its own repairs. These questions are
taken up, in order, by E2; by E1 and E3; by E9; by E6 and E8; and by E4 and E13. None of them is
answered in full here.

\paragraph{Biological autonomy.} Theories of autopoiesis, closure to efficient causation and
organizational closure describe living systems as maintaining the constraints that maintain them
\citep{MaturanaVarela1980,Rosen1991,MorenoMossio2015,MonteilMossio2015}. This paper does not
identify those theories with its definitions. It takes from them one distinction: the persistence
of an object is a different condition from the reinstatement of the apparatus that keeps it viable.
E3 states the two as separate certificates. Biological individuality \citep{Clarke2010,GodfreySmith2009},
trained immunity \citep{NeteaEtAl2016}, allostasis \citep{SterlingEyer1988} and bioelectric pattern
memory \citep{PezzuloLevin2015} appear only as readings that an application would have to support
with the data of \Cref{sec:scope}.

\paragraph{Regulation, optimization and memory.} Ashby's law of requisite variety and the good
regulator theorem are conceptual neighbours of E2 \citep{Ashby1956,ConantAshby1970}. E6 recovers
the Karush--Kuhn--Tucker conditions and the shadow-price reading of a multiplier
\citep{BoydVandenberghe2004}; rational inattention is a special case when discharge is measured by
mutual information \citep{Sims2003}. The firm as a boundary chosen to economize on transactions and
the Lucas critique are neighbours of E12 and E11 \citep{Coase1937,Lucas1976}. Reconsolidation after
retrieval and offline consolidation are the neuroscience reference points for E14 and E15
\citep{NaderSchafeLeDoux2000,Dudai2004,WalkerStickgold2006}. None of these results is re-proved.

\paragraph{Provenance of the results.} \Cref{tab:provenance} gives the origin of the results that draw on
earlier work, with the source and what is added here; the imported tools are listed in
\Cref{sec:intro-claims}.

\begingroup\small
\begin{longtable}{@{}>{\raggedright\arraybackslash}p{3.2cm}>{\raggedright\arraybackslash}p{3.2cm}>{\raggedright\arraybackslash}p{8.6cm}@{}}
\caption{Provenance of the results.}\label{tab:provenance}\\
\toprule
Result & Category & Source and what is new \\
\midrule
\endfirsthead
\toprule
Result & Category & Source and what is new \\
\midrule
\endhead
\bottomrule
\endfoot
\cref{prop:e2-self} & new; same-level self-soundness theorem imported & The same-level self-audit theorem of \FoundIII{} restricted to a carried level, with the weaker conclusion that the verdict is not \emph{accepted} \citep{Tsiokos2026FoundIII}. \\
\cref{thm:e6-kkt}, \cref{prop:e6-sufficient}, \cref{thm:e6-caps} & new; KKT conditions and concave sufficiency imported & \Cref{thm:e6-kkt} is the complementary-slackness reading of the Karush--Kuhn--Tucker conditions; the observation that a positive price needs a selected probe of positive discharge, and not only a binding budget, is the one point added here. \Cref{prop:e6-sufficient} is their standard sufficiency for concave problems \citep[\S 5.5.3]{BoydVandenberghe2004}; it is argued on paper and not formalized. \Cref{thm:e6-caps} is a separate sufficiency result, proved with cap multipliers, for capped linear objectives, which can fail to be differentiable at a cap, as in \Cref{thm:e6-atcap}. \\
\cref{thm:e9}, \cref{prop:e9-saturation} & new; saturation specialized & Existence of a maximum over a finite list is elementary; the saturation proposition is the same-family corollary of the adequacy-residual development specialized to the trace, with the chain rule as a hypothesis \citep{Tsiokos2026Adequacy}. The question it answers, which probe to add and at what price, is the probe-enlargement problem \citep{Tsiokos2026NeedleKiller}. \\
\end{longtable}
\endgroup
